\documentclass[12pt]{article}
\usepackage[utf8]{inputenc}
\usepackage[russian]{babel}
\usepackage{amsmath}
\usepackage{amssymb}
\usepackage{geometry}
\geometry{a4paper, margin=2cm}

\begin{document}

\section*{Задача 3.1.6e}

\textbf{Условие:} Известно, что $|a_i|\leq 100$. Реализовать сортировку подсчётом со сложностью $O(n)$ при равномерном весовом критерии.\\
\texttt{in:} $n, a_1, \ldots, a_n$\\
\texttt{out:} $b_1, \ldots, b_n$

\section*{Псевдокод}

\begin{verbatim}
def counting_sort(arr):
    min_val, max_val = -100, 100
    range_size = max_val - min_val + 1
    count = [0] * range_size

    for x in arr:
        count[x - min_val] += 1

    pos = 0
    for idx in range(range_size):
        value = idx + min_val
        for _ in range(count[idx]):
            arr[pos] = value
            pos += 1
    return arr
\end{verbatim}

\section*{Решение}

\begin{verbatim}
; Data Model
;0: r0
;1: trash | a_i | iterator
;2: n | a_i
;3 ... 203: arr
;204: [arr.end()]

Read 2
Load =0
Sub 2
JZero END
JGTZ END    ; if n <= 0 then exit

READ_DATA: Load 2
    JZero PRINT_INIT; if n == 0 then goto PRINT
    Read 1          ; r1 := read() = a_i
    Load 1
    Add =103        ; r0 := &arr[a_i]
    Store 1         ; r1 := &arr[a_i]
    Load *0         ; r0 := arr[a_i]
    Add =1
    Store *1        ; arr[a_i]++
    Load 2
    Sub =1
    Store 2         ; n--
    Jump READ_DATA
PRINT_INIT: Load =3
    Store 1         ; r1 := arr.begin()
PRINT: Load =204
    Sub 1
    JZero END       ; if r1 == arr.end() then exit PRINT
    Load 1
    Sub =103
    Store 2         ; r2 := a_i
    Load *1         ; r0 := arr[a_i]
; if arr[a_i] = k > 0 then print a_i k times
PRINT_NUM: JZero PRINT_AFTER_NUM    ; if r0 == 0 then end print a_i
    Write 2         ; print a_i
    Sub =1          ; r0--
    Jump PRINT_NUM
PRINT_AFTER_NUM: Load 1
    Add =1
    Store 1         ; r1++
    Jump PRINT
END: Halt 1
\end{verbatim}

\section*{Сложность}

При загрузке каждый входной элемент читается и увеличивает счётчик своего значения. Все $201$ ячеек счётчиков просматриваются в любом случае. Если все значения разные, то для каждого происходит один PRINT\_NUM. Если все данные одинаковые, то для одной ячейки происходит PRINT\_NUM. Но это не влияет на сложность, так как в PRINT\_NUM происходит только Write, а нужные данные сохраняются на этапе PRINT. Значит, для всех входов число шагов одинаково. Для удобства дальше рассмотрим случай, когда все входные значения различные.

\begin{enumerate}
    \item Проверка $n$ - $5$ шагов
    \item Чтение $n$ входных значений:
    \begin{enumerate}
        \item $13$ шагов на каждый элемент
    \end{enumerate}
    \item $2$ шага на завершающую проверку цикла
    \item PRINT\_INIT - $2$ шага
    \item Просмотр всех $201$ ячеек счётчиков и вывод:
    \begin{enumerate}
        \item Начало PRINT - 7 шагов
        \item PRINT\_NUM - 5 шагов для тех, что надо выводить (по предположению для удобства они все разные). 1 шаг для других
        \item PRINT\_AFTER\_NUM - 4 шага
    \end{enumerate}
    \item Финальная проверка условия цикла - 3 шага
    \item Halt 1 - 1 шаг
\end{enumerate}

Значит, сложность в худшем случае:
\[
    T_{\max}(n)=5 + 13*n + 2 + 2 + (201*(7 + 4) + 5*n + 1*(201 - n)) + 4= 17n + 2425
\]

\end{document}